\section{Controlled benchmark results}

\subsection{Aggregate accuracy}

The analysis includes the complete grid of \ConfirmatoryRuns{}
trained models and \ConfirmatoryTrajectories{} trajectory-level rows. All
prespecified metrics were finite (\AllMetricsFinite), and the 40 seed--system
groups produced \UniqueDataHashes{} distinct data fingerprints while sharing
one fingerprint across methods within each group.

Table~\ref{tab:main-results} reports the seed-balanced mean and standard
deviation for horizons 1, 10, and 25. Across all 11 methods, \BestMethod{}
had the lowest aggregate NRMSE-25 (\BestNRMSE). CFR-WM's mean was
\CFRNRMSE{}, compared with \PGNNNRMSE{} for the preselected PGNN-style comparator.

\begin{table*}[t]
\centering
\caption{Confirmatory accuracy over the fixed five-system, three-regime
benchmark. Entries are means $\pm$ SD across eight seed-level, cell-balanced
aggregates. Bold marks the lowest aggregate NRMSE-25. Parameter counts are means
over the five system-specific input dimensions; all methods received the same six-trial
validation search budget.}
\label{tab:main-results}
\resizebox{\textwidth}{!}{\begin{tabular}{lrrrr}
\toprule
Method & Mean parameters & NRMSE-1 & NRMSE-10 & NRMSE-25 \\
\midrule
CFR-WM & 57694 & 0.039 $\pm$ 0.001 & 0.057 $\pm$ 0.002 & \textbf{0.100 $\pm$ 0.004} \\
CFR-gain-0.25 & 57694 & 0.037 $\pm$ 0.001 & 0.059 $\pm$ 0.002 & 0.103 $\pm$ 0.004 \\
CFR-gain-0.75 & 57694 & 0.047 $\pm$ 0.001 & 0.065 $\pm$ 0.002 & 0.107 $\pm$ 0.005 \\
PGNN-style & 61275 & 0.060 $\pm$ 0.001 & 0.070 $\pm$ 0.002 & 0.110 $\pm$ 0.003 \\
CFR-gain-1 & 57694 & 0.059 $\pm$ 0.001 & 0.073 $\pm$ 0.002 & 0.114 $\pm$ 0.004 \\
PG-DMM-style & 61493 & 0.063 $\pm$ 0.001 & 0.075 $\pm$ 0.002 & 0.114 $\pm$ 0.003 \\
RSSM-style & 62485 & 0.073 $\pm$ 0.002 & 0.090 $\pm$ 0.005 & 0.123 $\pm$ 0.006 \\
Transformer-style & 59708 & 0.074 $\pm$ 0.002 & 0.094 $\pm$ 0.004 & 0.132 $\pm$ 0.009 \\
SimNorm-style & 62485 & 0.074 $\pm$ 0.002 & 0.098 $\pm$ 0.006 & 0.138 $\pm$ 0.013 \\
CFR-gain-0 & 57694 & 0.081 $\pm$ 0.003 & 0.100 $\pm$ 0.003 & 0.144 $\pm$ 0.005 \\
CFR-one-step & 57694 & 0.040 $\pm$ 0.001 & 0.136 $\pm$ 0.050 & 0.278 $\pm$ 0.104 \\
\bottomrule
\end{tabular}
}
\end{table*}

For the primary estimand, PGNN-style and CFR-WM
NRMSE-25 are first averaged with equal weight over the 15 prespecified
system--regime cells; their relative improvement is then computed. The
confirmatory point estimate is the arithmetic mean of those eight seed-level
relative improvements rather than the ratio of the two pooled means in
Table~\ref{tab:main-results}.

\subsection{Primary and gated robustness hypotheses}

Against PGNN-style, the seed-mean relative NRMSE-25 improvement was
\PrimaryImprovement\% (95\% paired fixed-cell bootstrap CI
[\PrimaryRelativeCILow\%, \PrimaryRelativeCIHigh\%]). In absolute units,
CFR-WM minus PGNN-style was \PrimaryDifference{} (95\% CI [\PrimaryCILow,
\PrimaryCIHigh]). The exact one-sided sign-flip test against the same 5\%
practical margin gave $p=\PrimaryPValue$. Under the prespecified conjunction of a
5\% interval lower bound, $p<0.05$, and finite metrics, H1 was
\PrimaryDecision.

For the H2 components, the combined corrupt-ID and corrupt-OOD regimes had a
relative effect of \CorruptImprovement\% (95\% CI
[\CorruptRelativeCILow\%, \CorruptRelativeCIHigh\%], one-sided 5\%-margin
$p=\CorruptPValue$). ID-clean had an effect of \CleanImprovement\% (95\% CI
[\CleanRelativeCILow\%, \CleanRelativeCIHigh\%], one-sided signed
$-5\%$-margin $p=\CleanPValue$). Because H2 is hierarchically gated by H1 and
requires both component rules, the robustness hypothesis was
\RobustnessDecision.

Table~\ref{tab:seed-effects} reports the eight independent cluster effects.
Every seed exceeded the 5\% practical threshold.

\begin{table}[t]
\centering
\caption{Seed-level primary effects after equal weighting of the 15
benchmark cells. Positive values favor CFR-WM.}
\label{tab:seed-effects}
\begin{tabular}{rr}
\toprule
Seed & CFR-WM improvement over PGNN-style, \% \\
\midrule
41 & +5.29 \\
53 & +12.76 \\
67 & +12.01 \\
79 & +10.49 \\
83 & +9.69 \\
97 & +6.66 \\
109 & +9.32 \\
113 & +5.56 \\
\bottomrule
\end{tabular}

\end{table}

\begin{figure*}[t]
\centering
\includegraphics[width=0.96\textwidth]{figures/confirmatory_overview.png}
\caption{Confirmatory summary. (A) Aggregate transition accuracy for the
six main methods. (B) Paired CFR-WM effects against PGNN-style by regime; positive
values favor CFR-WM. (C) Prespecified fixed-gain sensitivity. (D) Paired
fixed-cell intervals against each broad comparator. Error bars in A--C are
seed-level SD; D uses the nested 95\% bootstrap interval.}
\label{fig:confirmatory-overview}
\end{figure*}

\subsection{Comparator and sensitivity evidence}

Table~\ref{tab:comparator-contrasts} reports every prespecified broad
comparison. Positive relative effects favor CFR-WM. The confidence intervals
retain all fixed cells, whereas the exact two-sided seed sign-flip values are
Holm-adjusted as one five-test family. These secondary comparisons contextualize
the primary PGNN-style contrast.

The PGNN-style row repeats the primary point estimate but uses the
comparator-family bootstrap seed (20260808), whereas H1 uses its prespecified
primary seed (20260805). The estimand and 5,000-replicate resampling scheme are
otherwise identical; the separately sampled interval endpoints differ by at
most 0.01 percentage point after rounding, and only the primary interval enters
the H1 decision.

\begin{table*}[t]
\centering
\caption{CFR-WM relative NRMSE-25 effects against broad comparators. Exact
$p$ values test a zero seed-level difference; Holm adjustment is confined to
this prespecified family of five.}
\label{tab:comparator-contrasts}
\begin{tabular}{lrrr}
\toprule
Reference & Relative effect [95\% CI], \% & Exact $p$ & Holm $p$ \\
\midrule
RSSM-style & +19.08 [+17.16, +20.98] & 0.0078 & 0.0391 \\
SimNorm-style & +27.34 [+24.00, +30.62] & 0.0078 & 0.0391 \\
Transformer-style & +24.09 [+21.65, +26.64] & 0.0078 & 0.0391 \\
PGNN-style & +8.97 [+6.89, +10.96] & 0.0078 & 0.0391 \\
PG-DMM-style & +12.44 [+10.51, +14.39] & 0.0078 & 0.0391 \\
\bottomrule
\end{tabular}

\end{table*}

The lowest aggregate NRMSE-25 among gains $0$, $0.25$, $0.5$, $0.75$, and $1$
occurred at gain \BestGain. Relative to the one-step-trained variant, CFR-WM's
effect was \OneStepImprovement\% (95\% CI [\OneStepRelativeCILow\%,
\OneStepRelativeCIHigh\%], Holm $p=\OneStepHolmP$). Table
\ref{tab:sensitivity-contrasts} gives the complete second five-test family.

\begin{table*}[t]
\centering
\caption{CFR-WM relative NRMSE-25 effects against gain and training-horizon
sensitivities. Positive values favor the proposed gain-0.5 multi-step model.}
\label{tab:sensitivity-contrasts}
\begin{tabular}{lrrr}
\toprule
Reference & Relative effect [95\% CI], \% & Exact $p$ & Holm $p$ \\
\midrule
Gain 0 & +30.89 [+28.95, +32.67] & 0.0078 & 0.0391 \\
Gain 0.25 & +2.97 [+0.88, +5.27] & 0.0312 & 0.0391 \\
Gain 0.75 & +6.97 [+4.77, +9.42] & 0.0078 & 0.0391 \\
Gain 1 & +12.13 [+10.31, +13.92] & 0.0078 & 0.0391 \\
One-step training & +59.38 [+48.88, +68.88] & 0.0078 & 0.0391 \\
\bottomrule
\end{tabular}

\end{table*}

Table~\ref{tab:secondary} reports secondary diagnostics of Gaussian forecast
quality, marginal calibration, and catastrophic-rollout frequency.

\begin{table*}[t]
\centering
\caption{Secondary aggregate diagnostics for the six broad methods. Rates and
scores are seed-balanced over all fixed cells and are not separately tested.}
\label{tab:secondary}
\resizebox{\textwidth}{!}{\begin{tabular}{lrrrr}
\toprule
Method & Parameters & Gaussian NLL & Calibration error & Catastrophic rate \\
\midrule
RSSM-style & 62485 & -0.183 & 0.214 & 0.01\% \\
SimNorm-style & 62485 & -0.246 & 0.206 & 0.01\% \\
Transformer-style & 59708 & -0.279 & 0.202 & 0.01\% \\
PGNN-style & 61275 & -0.385 & 0.248 & 0.00\% \\
PG-DMM-style & 61493 & 0.029 & 0.245 & 0.00\% \\
CFR-WM & 57694 & 0.013 & 0.228 & 0.00\% \\
\bottomrule
\end{tabular}
}
\end{table*}

\subsection{Fixed-cell heterogeneity}

Figure~\ref{fig:cell-heterogeneity} retains all 15 system--regime effects. By
the prespecified descriptive guard, CFR-WM was more than 10\% worse than PGNN-style in
\CellsWorseTen{} cells. All three were ID-clean: cart--pole
(\CartpoleCleanChange\% relative change), controlled Van der Pol
(\VanDerPolCleanChange\%), and Duffing (\DuffingCleanChange\%). Conversely,
\IDCorruptFavoredCount{} of \IDCorruptCellCount{} ID-corrupt cells favored
CFR-WM, with descriptive changes from \IDCorruptLeastChange\% to
\IDCorruptMostChange\%. This pattern explains why an aggregate H1 success coexists with
the clean-regime noninferiority result. These cell-level estimates are
descriptive; their unadjusted bootstrap intervals are reported in the appendix
table and machine-readable CSV.

\begin{figure}[t]
\centering
\includegraphics[width=\linewidth]{figures/cell_difference_heatmap.png}
\caption{Descriptive CFR-WM change relative to PGNN-style in each fixed benchmark
cell. Negative values mean lower NRMSE-25 for CFR-WM.}
\label{fig:cell-heterogeneity}
\end{figure}

\subsection{Component analyses and stress tests}

Table~\ref{tab:posthoc-baseline} separates the context filter from the learned
residual. The gain-0.5 filter followed by nominal rollout reaches 0.111
NRMSE-25, below both CFR-WM (0.119) and PGNN-style (0.132). Evaluating the
trained CFR-WM weights at their stored residual multiplier of 0.25 reproduces the stored
cell means with maximum absolute difference
\ResidualReconstructionDifference{}. Across the four tested inference
multipliers, the lowest aggregate and regime-specific means occur at scale
\PosthocBestResidualScale{}. The learned residual therefore increases the
fixed-weight mean error in this analysis.

\begin{table*}[t]
\centering
\caption{Classical baselines and inference-time residual-scale
analysis. Entries are seed-balanced mean NRMSE-25 across the indicated fixed
cells. The scale analysis holds trained CFR-WM weights fixed and isolates the
residual contribution.}
\label{tab:posthoc-baseline}
\resizebox{\textwidth}{!}{\begin{tabular}{lrrrr}
\toprule
Analysis and variant & All & Clean & ID corrupt & OOD corrupt \\
\midrule
Nominal, final imputation & 0.137 & 0.040 & 0.151 & 0.221 \\
Complementary filter + nominal & 0.111 & 0.042 & 0.092 & 0.201 \\
PGNN-style & 0.132 & 0.052 & 0.134 & 0.211 \\
CFR-WM & 0.119 & 0.055 & 0.097 & 0.204 \\
\midrule
CFR-WM, inference scale 0 & 0.111 & 0.042 & 0.092 & 0.201 \\
CFR-WM, inference scale 0.125 & 0.113 & 0.046 & 0.093 & 0.202 \\
CFR-WM, inference scale 0.25 & 0.119 & 0.055 & 0.097 & 0.204 \\
CFR-WM, inference scale 0.5 & 0.135 & 0.079 & 0.113 & 0.213 \\
\bottomrule
\end{tabular}
}
\end{table*}

The synchronous two-step all-channel outage produces a 2.7\% increase over
clean-context CFR-WM, compared with 51.9\% for PGNN-style
(Table~\ref{tab:structured-outage}). For the pendulum, unit-circle projection
changes \ProjectionChangedFraction\% of the \ProjectionMissingEntries{}
missing sine/cosine entries after the direct correction. The mean, 95th
percentile, and maximum absolute normalized changes are \ProjectionMeanChange{},
\ProjectionPChange{}, and \ProjectionMaxChange{}, respectively. Thus the mask
sets the direct measurement gain to zero, but projection can still alter a
missing coordinate.

\begin{table*}[t]
\centering
\caption{Clean-ID performance after a synchronous two-step all-channel context
outage. Means $\pm$ SD and relative changes are across eight seed-level
aggregates after averaging systems and the three outage positions.}
\label{tab:structured-outage}
\begin{tabular}{lrr}
\toprule
Method & Outage NRMSE-25 & Change from clean (\%) \\
\midrule
PGNN-style & 0.079 $\pm$ 0.004 & +51.9 \\
CFR-WM & 0.057 $\pm$ 0.003 & +2.7 \\
\bottomrule
\end{tabular}

\end{table*}

Table~\ref{tab:latency} reports the stored measured runtime. CFR-WM requires
\CFRLatencyMean{} $\pm$ \CFRLatencySD{} ms per trajectory, including its
nominal context filtering. The measurement is a throughput statistic for one
96-trajectory GPU batch, not single-trajectory online latency.

\begin{table}[t]
\centering
\caption{Measured 25-step rollout time on the RTX 5070. Each cell uses one
batch of 96 trajectories and includes transfer, context encoding, and rollout;
SD is across seed-level cell-balanced means.}
\label{tab:latency}
\begin{tabular}{lrr}
\toprule
Method & ms/trajectory & SD \\
\midrule
RSSM-style & 0.110 & 0.006 \\
SimNorm-WM & 0.113 & 0.005 \\
Transformer-style & 0.109 & 0.007 \\
PGNN-style & 0.347 & 0.037 \\
PG-DMM-style & 0.344 & 0.036 \\
CFR-WM & 0.416 & 0.015 \\
\bottomrule
\end{tabular}

\end{table}
